%قبل از شروع مسئله ترجیحاً چیزی ننویسید. اما اگر خواستید بنویسید، در قالب لیستی (فقرات و شمارش) بنویسید.
%برای اضافه کردن نمره، بعدی از شروع‌مسئله، [x] را اضافه کنید و به جای x  نمره مورد نظر را قرار دهید. برای زیرمسئله‌ها هم به همین صورت است.
%\شروع‌مسئله[10]

\شروع‌مسئله[10]
یگ گراف با \متن‌لاتین{n}‌راس و \متن‌لاتین{m} یال داریم. الگوریتمی با \متن‌لاتین{O(n + m)} ارائه کنید که بگوید گراف دوبخشی است یا خیر. 
\\ گراف دوبخشی:‌ گرافی که بتوان رئوس آن را به دو دسته تقسیم کرد بطوریکه بین رئوس درون یک دسته یالی وجود نداشته باشد.

%اگر مسئله چند قسمت داشت به این صورت عمل کنید:
%\شروع{زیرمسائل}
%\زیرمسئله[0] % مسئلۀ ۱آ
%\زیرمسئله % مسئلۀ ۱ب
%\زیرمسئله % مسئلۀ ۱پ
%\زیرمسئله % مسئلۀ ۱ت
%\پایان{زیرمسائل}

\solution{
مساله معادل رنگ کردن راس‌های گراف با دو رنگ است به طوری که بین رئوس همرنگ یالی نباشد (رئوس هر رنگ را یک بخش از گراف دو بخشی در نظر می‌گیریم). برای این منظور، هر مولفه‌ی همبندی را بطور جداگانه بررسی می‌کنیم. اگر بتوان رئوس هر مولفه را به شکل گفته شده رنگ کرد (یعنی اگر هر مولفه از گراف دوبخشی باشد)، در این صورت کل گراف را هم می‌توان رنگ کرد چون رئوس هر مولفه به دو بخش تقسیم شده اند و بین رئوس دو مولفه‌ی همبندی جدا هم یالی نیست (وگرنه در بک مولفه‌ی همبندی بودند). اگر هم مولفه‌ای در گراف وجود داشته باشد که دوبخشی نباشد در این صورت کل گراف هم دوبخشی نیست. برهان خلف: فرض کنید مولفه‌ای دوبخشی نباشد اما کل گراف دوبخشی باشد: در اینصورت میتوان رئوس کل گراف را با دو رنگ رنگ‌آمیزی کرد که بین رئوس یک رنگ یالی نباشد و در نتیجه در همان رنگ‌آمیزی گراف، رئوس مولفه‌ی مدنظر هم به دو رنگ تقسیم می‌شوند که بین رئوسشان یالی نیست پس مولفه‌ی مدنظر هم دوبخشی بوده که با فرض اولیه ما (که دوبخشی نبودن این مولفه است) تناقض دارد.
\\حال برای چک کردن دوبخشی بودن هر مولفه، از یکی از رئوس آن مولفه \متن‌لاتین{dfs} می‌زنیم. به راسی که از آن \متن‌لاتین{dfs} زده‌ایم رنگ 1 را نسبت می‌دهیم (بنابه تقارن فرقی نمی‌‍کند به راس اول کدام رنگ را نسبت دهیم و می‌توان همه رنگ های را برعکس کرد تا در جواب رنگ این راس همان رنگ مدنظر ما باشد). حال به فرزندان این راس در زیردرخت \متن‌لاتین{dfs} مجبوریم رنگ 2 را نسبت بدهیم تا دو سر یال‌ها همرنگ نباشد. به همین شکل به هر راس در درخت \متن‌لاتین{dfs} مجبوریم رنگ مخالف با رنگ پدرش را نسبت بدهیم (یعنی از ریشه شروع می‌کنیم و به هر راس که در \متن‌لاتین{dfs} می‌رسیم رنگ مخالف با رنگ پدرش را نسبت می‌دهیم چون اگر رنگ پدرش مشخص باشد مجبوریم این رنگ را به این راس نسبت دهیم تا بین رئوس همرنگ یال نباشد) که این کار با \متن‌لاتین{o(n + m)} قابل انجام است. پس اگر رنگ ریشه را مشخص کنیم، رنگ تمام رئوسی که در \متن‌لاتین{dfs} به آن‌ها می‌رسیم به طور اجباری مشخص می‌شود و چون در \متن‌لاتین{dfs} به تمام رئوس مولفه می‌رسیم، پس مشخص می‌شود اگر گراف بخواهد دوبخشی باشد رنگ هر راس چه خواهد بود. بعد از آن که رنگ ها را نسبت دادیم برای هر یال بررسی می‌کنیم آیا دوسر آن همرنگ هستند یا خیر (با \متن‌لاتین{o(m)}) و اگر یالی با دوسر همرنگ پیدا شود گراف دوبخشی نیست (چرا که رنگ‌های رئوس به طور یکتا مشخص شده‌اند) و اگر هم چنین یالی پیدا نشد گراف دوبخشی بوده و با رنگ‌آمیزی ارائه شده آن‌ها را به دوبخش تقسیم کرده‌ایم. پس با \متن‌لاتین{o(n + m)} جواب را بدست می‌آوریم.
\\
}




%%%%%%%%%%%%%%%%%%%%%%%%%%%%%%%%%%%%%%%%%%%%%%%%%%%%%
% برخی دستورات مفید برای تمرین‌ها. %
%%%%%%%%%%%%%%%%%%%%%%%%%%%%%%%%%%%%%%%%%%%%%%%%%%%%%

% فرمول‌های نمونه:
%$f(x) = \Theta(x)$
%$T(x, y) \leq \log(x) + 2^y + \binom{2n}{n}$
% {\tt code\_function}


% لیست موارد بدون شماره‌گذاری:
%\شروع{فقرات}
%    \فقره مورد اول 
%        \شروع{فقرات}
%            \فقره مورد اول داخلی 
%                \شروع{فقرات}
%                    \فقره اولین مورد خیلی داخلی 
%                    \فقره دومین مورد خیلی داخلی 
%                \پایان{فقرات}
%            \فقره مورد دوم داخلی 
%        \پایان{فقرات}
%    \فقره مورد دوم 
%\پایان{فقرات}

% لیست موارد شماره‌گذاری شده:
%\شروع{شمارش}
%    \فقره یک مورد 
%    \فقره یک مورد دیگر 
%    \فقره باز هم یک مورد دیگر
%\پایان{شمارش}

% برای نوشتن یک فرمول چند خطی که یک شماره بخوره. برای حالت‌های دیگر به سایت زیر مراجعه کنید:
% https://www.overleaf.com/learn/latex/Aligning_equations_with_amsmath
%\begin{align} 
%    \begin{split}
%        (x+y)^3 &= (x+y)^2(x+y) \\
%                &= (x^2+2xy+y^2)(x+y) \\
%                &= (x^3+2x^2y+xy^2) + (x^2y+2xy^2+y^3) \\
%                &= x^3+3x^2y+3xy^2+y^3
%    \end{split}                                 
%\end{align}

% نوشتن ماتریس و آرایه دوبعدی
%\begin{align}
%    A = 
%    \begin{bmatrix}
%        A_{11} & A_{21} \\
%        A_{21} & A_{22}
%    \end{bmatrix}
%\end{align}

% اضافه کردن تصویر با شرح. برای اضافه کردن تصویر بدون شرح، فقط دستور درج‌تصویر را استفاده کنید.
%\شروع{شکل}[h!]
%\تنظیم‌ازوسط
%  \درج‌تصویر[پهنا=0.5\پهنای‌متن]{img.jpg}
%\شرح{شرح تصویر}
%\برچسب{تصویر:۱}
%\پایان{شکل}
%\پایان{زیرمسائل}

%کد زیر را ببینید.
%%\lstinputlisting{source_filename.py}
%\begin{latin}
%\begin{lstlisting}
%# Import the modules
%import sys
%import random
%
%ans = True
%
%while ans:
%question = raw_input("Ask the magic 8 ball a question: (press enter to quit) ")
%
%answers = random.randint(1,8)
%
%if question == "":
%sys.exit()
%
%elif answers == 1:
%print "It is certain"
%
%elif answers == 2:
%print "Outlook good"
%
%elif answers == 3:
%print "You may rely on it"
%
%elif answers == 4:
%print "Ask again later"
%
%elif answers == 5:
%print "Concentrate and ask again"
%
%elif answers == 6:
%print "Reply hazy, try again"
%
%elif answers == 7:
%print "My reply is no"
%
%elif answers == 8:
%print "My sources say no"
%\end{lstlisting}
%\end{latin}